\documentclass[11pt]{article}
\usepackage[a4paper,margin=21mm]{geometry}
\usepackage[no-math]{fontspec}
\setmainfont{Latin Modern Roman}
\newfontfamily\CJKfont{Noto Serif CJK SC}
\XeTeXlinebreaklocale "zh"
\XeTeXlinebreakskip=0pt plus 1pt
\usepackage{amsmath,amssymb,amsthm,amscd,graphicx,color}
\usepackage{xurl}
\usepackage[unicode,hidelinks]{hyperref}
\allowdisplaybreaks
\setlength\emergencystretch{3em}
\numberwithin{equation}{section}

\usepackage[all]{xy}
\usepackage{multirow}

\numberwithin{equation}{section}

\newtheorem{Theorem}{Theorem}[section]
\newtheorem{Corollary}[Theorem]{Corollary}
\newtheorem{Lemma}[Theorem]{Lemma}
\newtheorem{Proposition}[Theorem]{Proposition}
\newtheorem{Conjecture}[Theorem]{Conjecture}
\newtheorem{theoremx}[Theorem]{Theorem}

\usepackage{mathrsfs}
\DeclareMathAlphabet{\mathscrbf}{OMS}{mdugm}{b}{n}

\usepackage[x11names]{xcolor}

\definecolor{burgundy}{rgb}{0.5, 0.0, 0.13}
\definecolor{carmine}{rgb}{0.59, 0.0, 0.09}
\definecolor{bostonuniversityred}{rgb}{0.8, 0.0, 0.0}
\definecolor{alizarin}{rgb}{0.82, 0.1, 0.26}

\def\burgundy#1{\textcolor{bostonuniversityred}{#1}}

\newcommand{\changed}[1]{{\color{bostonuniversityred} #1}}
% \newcommand{\changed}[1]{{#1}}

\newcommand{\rmk}[1]{{\color{blue} {\rm [#1]}}}
\def\SG#1{\textcolor[rgb]{0.01,0.75,0.24}{[#1]}}

\def\BN{\mathbb N}
\def\BZ{\mathbb Z}
\def\BP{\mathbb P}
\def\BQ{\mathbb Q}
\def\BR{\mathbb R}
\def\BC{\mathbb C}
\def\BE{\mathbb E}
\def\BT{\mathbb T}
\def\BK{\mathbb K}
\def\BJ{\mathbb J}

\def\calA{\mathcal A}
\def\calB{\mathcal B}
\def\calC{\mathcal C}
\def\calD{\mathcal D}
%\def\calK{\mathcal K}
\def\calL{\mathcal L}
\def\calG{\mathcal G}
\def\calT{\mathcal T}
\def\calM{\mathcal M}
\def\calN{\mathcal N}
\def\calF{\mathcal F}
\def\calV{\mathcal V}
\def\calW{\mathcal W}
\def\calP{\mathcal P}
\def\calR{\mathcal R}
\def\calT{\mathcal T}
\def\calS{\mathcal S}
\def\calI{\mathcal I}
\def\calH{\mathcal H}
\def\calP{\mathcal P}
\def\calO{\mathcal O}
\def\calX{\mathcal X}
\def\calE{\mathcal E}

\def\D{\Delta}
\def\aa{\alpha}
\def\bb{\beta}
\def\gg{\gamma}

\def\La{\Lambda}
\def\l{\lambda}
\def\k{\kappa}
\def\w{\omega}

\def\la{\langle}
\def\ra{\rangle}

\def\ep{\epsilon}

\def\ti{\widetilde}
\def\wt{\widetilde}
\def\wh{\widehat}

\def\longto{\longrightarrow}
\def\SL{\mathrm{SL}}
\def\PSL{\mathrm{PSL}}
\def\GL{\mathrm{GL}}
\def\Sp{\mathrm{Sp}}
\def\pt{\partial}
\def\Ker{\mathrm{Ker}}

\def\z{\zeta}
\def\CS{\mathrm{CS}}
\def\Vol{\mathrm{Vol}}
\def\Om{\Omega}
%\def\Li{\mathrm{Li}}
\def\calK{\mathcal K}

\def\Phih{\widehat\Phi}
\def\Psih{\widehat\Psi}
%\def\Phio{\stackrel{\circ}{\varphi}}
\def\Phio{\widehat{\Phi}}
\def\phih{\widehat\phi}
\def\Psih{\widehat\Psi}



%%%%%%%%%%%%% Don's definitions
\def\e{\mathbf{e}}
\def\a{\alpha}
\def\p{\mathfrak p}
\def\l{\lambda}
\def\h{\hbar}
\def\={\;=\;}
\def\+{ + }

\def\R{\mathbb R}
\def\C{\mathbb C}
\def\Z{\mathbb Z}
\def\N{\mathbb N}
\def\Q{\mathbb Q}
\def\={\;=\;}
\def\+{ + }
\def\m{ - }
\def\si{\;\sim\;}
\def\i{^{-1}}
\def\e{\mathbf{e}}
\def\hb{\hbar}
\def\a{\alpha}
\def\b{\beta}
\def\l{\lambda}
\def\z{\zeta}
\def\G{\Gamma}
\def\s{\sigma}
\def\n{\nu}
\def\P{\Phi}
%\def\p{\varphi}
\def\v{\varepsilon}
\def\W{\kappa} % the weight, previously \l (=\lambda) sub-\s

\def\sma#1#2#3#4{\bigl(\smallmatrix#1&#2\\#3&#4\endsmallmatrix\bigr)} % small matrices
\def\vsma#1#2#3#4{(\smallmatrix#1&#2\\#3&#4\endsmallmatrix)} % very small matrices
\def\mat#1#2#3#4{\Bigl(\begin{matrix}#1&#2\\#3&#4\end{matrix}\Bigr)}
\def\L{\text{{\rm Li}}}
\def\den{\mathrm{den}}
\def\num{\text{num}}
\def\O{O}
\def\o{o}

\def\red#1{\textcolor{red}{#1}}
\def\blue#1{\textcolor{blue}{#1}}
\def\green#1{\textcolor[rgb]{0,0.0,1.0}{#1}}

\def\CH{{\color{red} \bf [CHECK!] }}
\def\G{\Gamma}
\def\J{\mathbf J} \def\J{J} % the J-function as an element on Q/Z
\def\Jt{\mathbf{\dot{\J}}}
\def\calJ{\mathcal J} % the J-function as an element of the Habiro ring
\def\BJ{{\mathbf Q}} % the matrix-valued version of this
\def\Qbar{\overline{\Q}}
\def\V{\mathbf V}
\def\tV{v} % = (cx'ed volume)/(2 pi i); was previously \def\tV{\widetilde V}
\def\be{\begin{equation}}
\def\ee{\end{equation}}
\def\ve{\varepsilon}
\def\ssm{\smallsetminus}

\def\K#1{{K_{#1}}}
\def\PP{\mathcal P}
\def\HH{\mathbb H}

\setcounter{secnumdepth}{5}
\setcounter{tocdepth}{5}
\def\Li{\mathrm{Li}}
\newcommand{\mb}{\mathbf}
\def\tcdot{\!\cdot\!}
\def\tiA{\wt A}
\def\tiB{\tilde{B}}
\def\opt{{\mathrm{opt}}}
\def\smo{{\mathrm{smooth}}}

\def\ul{\underline}
\def\g{\gamma}
\def\cE{\mathcal E}
\def\A{\mathbf A}
\def\bH{\mathbf H}
\def\B{\mathcal B} \def\hB{\widehat\B}
\def\bK{\mathbf K}
\def\bx{\mathbf x}
\def\bQ{Q}
\def\sQ{\mathcal{Q}}
%\def\diag{\operatorname{diag}} % (this is what Stavros had)
\def\diag{\operatorname{\mathbf{diag}}} % (boldface version)
\def\bJ{J} \def\bJ{\mathbf{J}}
\def\sJ{\mathcal{J}}
\def\pert{\mathrm{pert}}
\def\hol{\mathrm{hol}}
\def\conn{^\mathrm{conn}}
\def\sH{\mathcal H} % This "H" should be a "script H" since it is a
\def\th{\theta}
\def\Tr{\mathrm{Tr}}
\def\DD{DD}
\def\F{\mathcal F}
\def\bPhi{\mathbf\Phi}
\def\bPhih{\mathbf{\Phih}}
%\def\bphi{\mathbf\phi}
%\def\bphih{\mathbf{\phih}}
\def\jp{\mathbf{j}}
\def\WW{\Lambda^2}
\def\ww{ \wt\wedge }

\def\vve{E}
\def\jt{\jp_\l} % tweaked function with \lambda
\def\jt{\jp_{\text{tw}}} % tweaked function with "tw"
\def\jt{\jp_{\textsf T}} % tweaked function with "T", still ugly; suggestions??
\def\jt{\jp^{(\l)}}
\def\jt{\jp^{(\l)}}
\def\jt{\wt{\jp}}
\def\jhol{{\mathbf{j}}^{\mathrm{hol}}}
\def\EXACT{^{\text{\rm exact}}}
\def\RED{^{\text{red}}}

\def\bF{\mathbf F}
\def\FF{\mathfrak F}

\def\jpx{j\RED} % \def\jpx{j^*}
\def\bJx{\bJ\RED} % \def\bJx{\bJ^*}
\def\bPhihx{\widehat{\bPhi}\RED} \def\Phihx{\bPhihx}
\def\bPhix{\bPhi\RED} \def\Phix{\bPhix}
\def\Wx{W\RED}
\def\calPx{\mathcal{P}\RED} % \def\calPx{\mathcal{P}^*}

%\def\bphi{\mathbf\phi}
\def\bPsi{\BJ}
\def\MM{\pmb\mu}


\makeatletter\newlabel{K1}{{1.1}{}{Original source locator}{}{}}\makeatother
\makeatletter\newlabel{sec.perturbation}{{6}{}{Original source locator}{}{}}\makeatother
\makeatletter\newlabel{sec.Matrix}{{5}{}{Original source locator}{}{}}\makeatother
\makeatletter\newlabel{sub.PhiPerturbative}{{6.3}{}{Original source locator}{}{}}\makeatother
\makeatletter\newlabel{sub.CoeffA}{{3.4}{}{Original source locator}{}{}}\makeatother
\makeatletter\newlabel{cunit}{{9.3}{}{Original source locator}{}{}}\makeatother
\makeatletter\newlabel{Nahm3}{{6.7}{}{Original source locator}{}{}}\makeatother
\makeatletter\newlabel{newpsixz}{{6.8}{}{Original source locator}{}{}}\makeatother
\begin{document}
\begin{center}\Large Universal coefficient denominator bound for alpha-zero knot perturbative series defined by formal Gaussian integrals at nontrivial representations\end{center}
\noindent\textbf{Stable ID:} 989\quad\textbf{Author:} Stavros Garoufalidis; Don Zagier\par
\noindent\textbf{Work:} Knots, Perturbative Series and Quantum Modularity\par
\noindent\textbf{Source:} \url{https://sigma-journal.com/2024/055/}\par
\noindent\textbf{Version:} Official SIGMA 20 (2024) article 055, DOI 10.3842/SIGMA.2024.055; journal PDF and native TeX acquired 2026-09-30, exact bytes pinned by SHA256\par
\noindent\textbf{Rights:} CC BY 4.0. Authors retain copyright. \url{https://creativecommons.org/licenses/by/4.0/}. Reuse and typesetting adaptation; original language and mathematical content preserved.\par
\noindent\textbf{Difficulty (prior selection preserved):} {\CJKfont H3: The author bounds denominators of Gaussian averages through even symmetric exponent matrices and sharp prime-adic factorial floors. A separate decomposition of Bernoulli-weighted monomials controls the interaction between Gaussian pairing denominators and exponential multiplicities. The resulting universal D\_n bound is proved by a nontrivial prime-by-prime argument, including the complete Bernoulli denominator lemma.}\par
\medskip\noindent\textbf{Editorial boundary note (not author text):} Selected scope is exactly the previously reviewed alpha=0, c=1 nontrivial-representation formal Gaussian-series case of Theorem9.1. The original broader theorem is retained literally, together with the author explicit caveats that the proof is only for alpha=0, does not give the general rational-alpha details and excludes the trivial representation; those unproved extensions are not counted as this outcome. The proof adds the prime2 to the permitted finite excluded-prime set S, as the author expressly states. Full indexing/trace-field and normalized-coefficient definitions, universal denominator, exact Gaussian exponential/Laplacian definition, symmetric exponent-matrix bound, all four Bernoulli monomial cases, final prime-adic floor argument and entire deferred Bernoulli lemma are included. Original series-triangulation definitions and classical von Staudt–Clausen/factorial/multinomial facts remain precise references; quantum modularity, broader analytic asymptotics and general rational-alpha construction are not assumed proved here. The later connected-denominator theorem explicitly has an omitted proof and is excluded. Literal author formulas, inequalities, cases and omission caveats remain unchanged; no mathematical refereeing or editor proof.\par
\noindent\textbf{External original reference K1:} Original quantum-modularity asymptotic convention, source764–771, PDF9. Optional volume-normalization reference only; no QMC truth is assumed by this selected arithmetic proof.\par
\noindent\textbf{External original reference sec.perturbation:} Original half-symplectic/Neumann–Zagier background section, source3193–3778, PDF43–51. Formal earlier construction of the fixed algebraic shape data, referenced as background.\par
\noindent\textbf{External original reference sec.Matrix:} Original matrix-cocycle discussion, source2325–3191, PDF31–43. Unselected future discussion from indexing-set introduction.\par
\noindent\textbf{External original reference sub.PhiPerturbative:} Original perturbative-construction subsection, source3628–3778, PDF49–51. The selected proof recalls its exact Gaussian-series definition and Laplacian averaging at4864–4923; its formal knot-triangulation datum is an earlier cited definition, not a deferred current denominator argument.\par
\noindent\textbf{External original reference sub.CoeffA:} Original coefficient asymptotics subsection, source1527–1708, PDF20–22. Reference only identifies the alternative coefficient notation; no unproved asymptotic hypothesis is used in the selected bound.\par
\noindent\textbf{External original reference cunit:} Original near-unit Galois convention, source4717–4724, PDF64. Describes the prefactor convention, not an assumption of the selected alpha-zero Gaussian denominator proof.\par
\noindent\textbf{External original reference Nahm3:} Original Nahm-sum formula, source3707–3713, PDF50. Background relation to the quadratic form; exact current Gaussian exponential and averaging definition are retained in the selected proof.\par
\noindent\textbf{External original reference newpsixz:} Original Bernoulli-polynomial local series, source3732–3740, PDF51. Background normalization comparison, with the exact Bernoulli-number series used in this proof written at4867–4874.\par
\medskip\hrule\medskip\noindent\textbf{Author text begins below.}\par
\setcounter{section}{1}
% Original source lines 918--968
\section{A collection of formal power series}
\label{sec.Phi}

\subsection[The indexing set P\_K]{The indexing set $\boldsymbol{\calP_K}$}
\label{sub.p}

The power series mentioned above are labeled by a finite set $\calP_K$
that coincides with the set of boundary parabolic $\SL_2(\C)$-representations
of $\G:=\pi_1\big(S^3\ssm K\big)$ (or equivalently, of flat connections
on~$S^3\ssm K$ whose restriction to the peripheral subgroup of $\G$ are
parabolic) whenever the latter is finite. For a hyperbolic knot, this
set has three distinguished elements, denoted $\s_0$, $\s_1$ and $\s_2$,
corresponding respectively to the trivial representation, the geometric
representation (given by the natural embedding of~$\G$ into the isometry group
of~$\mathbb H^3$) and the antigeometric representation, which is its
complex conjugate and corresponds to the geometric representation of the
orientation-reversed hyperbolic knot. We denote by $\calP\RED_K=\calP_K\ssm\{\s_0\}$
the reduced set of non-trivial representations (or connections), and often
number the elements of $\calP_K$ as $\s_0,\s_1,\dots,\s_r$, where
$r:=\big|\calP\RED_K\big|$ will be called the \emph{rank} of the knot. (We hope that the
superscript ``red'' will not confuse the reader into thinking that the representations
in $\calP\RED_K$ are reducible; in fact, quite the opposite is true.)
The points of~$\calP\RED_K$ correspond to the solutions (in~$\C$) of a set of polynomial
equations \big(the so-called Neumann--Zagier equations coming from a triangulation
of~$S^3\ssm K$\big) as explained in detail in Section~\ref{sec.perturbation}.
In particular, $\calP\RED_K$ comes equipped with an action of the absolute
Galois group $\text{Gal}\big(\overline{\BQ}/\BQ\big)$, so to every element
$\s \in \calP\RED_K$ is associated a number field~$F_\s$ \big(given either as the
field generated by the coordinates of the solution of the NZ equations or
as the fixed field of the stabilizer of $\s$ in $\text{Gal}\big(\Qbar/\BQ\big)$\big),
called its trace field, together with an embedding, also denoted~$\s$,
of~$F_\s$ into $\Qbar\subset\C$. The field $F_{\s_1}$
coincides with the trace field~$F_K$ of~$K$ as introduced above and
$F_{\s_2}$ is the same field with the complex conjugate embedding into~$\C$.
Two more important invariants of $\s\in\calP_K\RED$ are an element $\xi_\s$
of the Bloch group (or third $K$-group) of~$F_\s$ and a complexified
volume $\V(\s)=\V(K,\s)$ (= ${\rm i}$ times the usual volume plus the Chern--Simons
invariant), obtained as the image of~$\xi_\s$ under the Borel regulator
map or via the dilogarithm, or alternatively its renormalized version
$v(\s)=v(K,\s)=\V(\s)/2\pi {\rm i}$, which for the geometric representation is
the same as the number $v(K)$ occurring in~\eqref{K1}.
(We will usually omit the letter~$K$ in this and all similar notations when
the knot is not varying.) We extend all of these invariants to $\calP_K$
by setting~$F_{\s_0}=\Q$, $\V(s_0)=0$, $\xi_{\s_0}=0$.
In Section~\ref{sec.perturbation} of Part~II,
we will give more details about the set $\calP_K$ and its invariants, and
also say something about the situation when the variety of parabolic
representations contains positive-dimensional components. In
Sections~\ref{sec.Matrix} and~\ref{sec.perturbation}, we will also describe
a large (matrix- rather than vector-valued) collection of power series
associated to~$K$.

\setcounter{section}{9}\setcounter{subsection}{1}\setcounter{equation}{1}
% Original source lines 4701--4708
\be
\label{Dn}
D_n = 2^{3n + v_2(n!)} \prod_{\substack{\text{$p$ prime} \\ p>2}}
 p^{ \sum_{i\ge0}[n/p^i(p-2)]} .
\ee
(Note that the exponent $v_p(D_n)$ of $p >2$ in $D_n$ can be written
as $r+v_p(r!) = v_p((p r)!)$ where $r=[n/(p-2)]$.) We will
return to this point in the next subsection.

\setcounter{equation}{3}
% Original source lines 4726--4748

We conjecture that these properties (a)--(g) hold for all hyperbolic knots~$K$, all
representations~$\s$ in~$\calP_K$ and all roots of unity $\z_\a$, with $F$ replaced
by the trace field~$F_\s$ and 5~by the denominator of~$\a$, as well as a~few other
small modifications (in particular, that instead of a~unit one may get an $S$-unit
for small finite set~$S$ of primes, essentially the ones occurring in the shape
parameters of a triangulation of $S^3\ssm K$, which was empty for $4_1$). In other
words, the power series \smash{$\Phi^{(K,\s)}_{\a}(h)$} can be written in the form
\be
\label{Phiprec2}
\Phi_\a^{(K,\s)}(h) =\mu_{\s,\a}\cdot(\ve_{\s,\a})^{1/c}
\cdot \delta_\s^{-1/2} \sum_{n=0}^\infty \tiA^{(K,\s)}_{\a,n} h^n ,
\qquad \tiA^{(K,\s)}_{\a,n} \in F_\s(\z_\a)
\ee
(so that \smash{$\tiA^{(K,\s)}_{\a,n}$} is the product of an algebraic number independent of~$n$
and the coefficient denoted \smash{$A^{(K,\s)}_\a(n)$} in Section~\ref{sub.CoeffA}),
where $F_\s$ is defined as in Section~\ref{sec.Phi}, $\mu_{\s,\a}$ is an $(8c)$-th root
of unity, and $\ve_{\s,\a} \in F_\s(\z_\a)^\times$ is a near-unit, canonically defined only
up to $c$-th powers, that transforms up to $c$-th powers as in~\eqref{cunit} (with 5
replaced everywhere by~$c$) and that conjecturally depends only on the element of the
Bloch group $B(F_\s)$ determined by $\s$ and in fact coincides with the near-unit that
was constructed in~\cite{CGZ}, and with the same denominator bound~$D_n$ as
in~\eqref{Dn}, independent of~$K$, $\s$ and~$\a$.


% Original source lines 4750--5097
\subsection{Denominators and integrality properties}
\label{sub.denominators}


The universal denominator statement given in formula~\eqref{Dn} above was found empirically
on the basis of the extensive numerical data for the $4_1$, $5_2$ and the $(-2,3,7)$ pretzel
knots presented in the appendix to this paper. In this section, we prove it for the denominators
of the power series defined in terms of Gaussian-type integrals in~\cite{DG2}.
This proof only applies to $\s\in\calP\RED_K$, since there is no such integral representation
for the trivial representation, but the corresponding denominator statement is true here also
and can in fact be strengthened because the power series in that case come
from the Habiro ring, as explained at the end of this section.

\begin{Theorem}
\label{thm.denom}
For each knot $K$, representation $\s\in\calP_K$, and number $\a \in \BQ/\BZ$, we have
\[%\label{AD}
D_n  \tiA^{(K,\s)}_{\a,n} \in \calO_S\big[\z_\a,c^{-1}\big],
\]
where \smash{$\Phi_\a^{(K,\s)}(h)$} is as in~{\rm \cite{DG2}}, $D_n$ is as in~\eqref{Dn},
\smash{$\tiA^{(K,\s)}_{\a,n}$} is as in~\eqref{Phiprec2}, $c$ is the denominator of~$\a$,~$\calO$ is the ring of integers of $F_\s$ and $S$ is a finite set of primes
of $F_\s$ that depends on $K$ but not on $n$ or on~$\a$.
\end{Theorem}

The first few values of $D_n$ are given by
\begin{align*}%\label{Dnvalues}
&   1, \ 24, \ 1152, \ 414720, \ 39813120,  \ 6688604160, \ 4815794995200, \  115579079884800, \\
&  22191183337881600, \ 263631258054033408000, \ 88580102706155225088000, \\
&   27636992044320430227456000, \ 39797268543821419527536640000, \ \dots.
\end{align*}
Campbell Wheeler pointed out to us that the above sequence appears (with no proof)
to equal to the sequence~\texttt{A144618} of the online-encyclopedia of integer
sequences~\cite{OEIS}, the latter related to Stirling's formula with half-shift
$D_n=\text{den}(a_n)$, where
\[%\label{stirling}
z! \sim \sqrt{2\pi} (z+1/2)^{z+1/2} {\rm e}^{-z-1/2}
\sum_{n=0}^\infty \frac{a_n}{(z+1/2)^n},
\qquad z \to \infty .
\]

The numbers $D_n$ grow rapidly, for example
\[
D_{50} = 2^{197} 3^{72} 5^{19} 7^{11} 11^5 13^4 17^3 19^2 23^2 29^1
 31^1 37^1 41^1 43^1 47^1
\]
or $D_{50}/24^{50}50!=5^7 7^3 11^1 13^1 17^1$. \big(In general, $D_n/24^nn!$ is an
integer whose $n$th root tends to~\smash{$\prod_{p\ge5}p^{2/(p-1)(p-2)}=1.8592481285\dots$}.\big)
We give the first 50 values of~$D_n$ by tabulating the ratio $\delta_n=D_n/3D_{n-1}$
(after removing the power of~2, and omitting the values equal to~1):

\begin{center}
\small
\begin{tabular}{|ll|ll|ll|ll|ll|ll|}
\hline
$n$ & $\delta_n$ & $n$ & $\delta_n$ & $n$ & $\delta_n$ & $n$ & $\delta_n$ &
$n$ & $\delta_n$ & $n$ & $\delta_n$
\\ \hline
3 & $3 \tcdot 5$ &
11 & $13$ &
20 & $7$ &
27 & $3^3 \tcdot 5 \tcdot 11 \tcdot 29\!$ &
35 & $7^2 \tcdot 37$ &
42 & $3 \tcdot 5 \tcdot 23$
\\ %\hline
5 & $7$ &
12 & $3 \tcdot 5$ &
21 & $3 \tcdot 5 \tcdot 23\!$ &
29 & $31$ &
36 & $3^2 \tcdot 5 \tcdot 11$ &
44 & $13$
\\ %\hline
6 & $3 \tcdot 5$ &
15 & $3 \tcdot 5^2 \tcdot 7 \tcdot 17\!$ &
22 & $13$ &
30 & $3 \tcdot 5^2 \tcdot 7 \tcdot 17$ &
39 & $3 \tcdot 5 \tcdot 41$ &
45 & $3^2 \tcdot 5^2 \tcdot 7 \tcdot 11 \tcdot 17 \tcdot 47\!$
\\ %\hline
9 & $3^2 \tcdot 5 \tcdot 11\!$ &
17 & $19$ &
24 & $3 \tcdot 5$ &
33 & $3 \tcdot 5 \tcdot 13$ &
40 & $7$ &
48 & $3 \tcdot 5$
\\ %\hline
10 & $7$ &
18 & $3^2 \tcdot 5 \tcdot 11$ &
25 & $7$ &
34 & $19$ &
41 & $43$ &
50 & $7$
\\ \hline
\end{tabular}

\end{center}

\noindent
We also remark that $D_{n_1}D_{n_2}|D_{n_1+n_2}$ for all $n_1,n_2\ge0$ and hence
that the subgroup
\[%\label{defRD}
R_D[[h]] = \left\{ \sum_{n=0}^\infty \frac{a_n}{D_n} h^n  \,\big|\,
a_n \in R \  \text{for all} \   n \right\}
\]
of $K[[h]]$ is a subring for every subring $R$ of a field~$K$ of characteristic zero.

\begin{proof}
We give the proof only for the case~$\a=0$, $c=1$, using the formulas in~\cite{DG}.
The general case can be proved along the same lines using the more complicated
formulas in~\cite{DG2}, in which the Bernoulli numbers are replaced by Bernoulli
polynomials, but we do not give the details here.
We will also ignore the prime~2 in our proof since it behaves somewhat differently
and in any case can be added to the finite set of excluded primes~$S$ in the statement
of the theorem.

The power series \smash{$\Phi_0^{(K,\s)}(h)$} attached to a triangulation $\calT$
of~$\mathbb S^3\ssm K$ were defined in~\cite{DG} as formal Gaussian integrals
$\langle f_{\calT}\rangle$
of the formal power series
\be
\label{psixz}
f_{\calT}(x;z) = \exp\Bigg(\sum_{j=1}^N \sum_{\substack{r, k\ge0 \\ 2r+k-2 > 0}}
\frac{B_r}{r!} \frac{(-x_j)^k}{k!} \Li_{2-r-k}(z_j)  h^{r+\tfrac k2-1} \Bigg)
\ee
in a multi-variable $x=(x_1,\dots,x_N)$, where $z_1,\dots,z_N$ are the shape parameters
of~$\calT$ and where~${\langle f\rangle=\langle f(x)\rangle_Q}$ is the mean value
defined by Gaussian integration with respect to a certain quadratic form~$Q$ with
coefficients in the field~$F_\s$. This form is essentially the one given by the
symmetric matrix $A+B \,\diag(z_j/(1-z_j))$ that occurred in the discussion
of~\eqref{Nahm3}
in Section~\ref{sub.PhiPerturbative}, and the function~\eqref{psixz} is essentially
the product of the functions occurring in~\eqref{newpsixz}, except that the terms
with $r+k=m$ fixed were combined there into a single Bernoulli polynomial $B_m(x)$
for~$m\ge3$ or $B_2(x)-x^2$ for~$m=2$, and that the normalizations used in~\cite{DG}
were slightly different from the ones used in Section~\ref{sub.PhiPerturbative}.

We now expand the exponential in~\eqref{psixz} as the product of the exponentials
of the monomials in the sum, and recall that $\Li_{2-m}(z)\in\Z[1/(1-z)]$ for
every~$m\ge2$, to deduce that \smash{$\Phi_\a^{(K,\s)}(h)=\langle f_{\calT}\rangle$} is an
$R$-linear combination of Gaussian averages $\langle T\rangle$ of terms~$T$ of the
form
\be
\label{term}
T = \prod_{j=1}^N\prod_{\substack{r, k\ge0 \\ 2r+k\ge3\substack}}
\frac1{\l_j(r,k)!} \biggl(\frac{B_r}{r!}
\frac{x_j^k}{k!} h^{r+\tfrac k2-1}\biggr)^{\l_j(r,k)}
\ee
with non-negative multiplicities $\l_j(r,k)$, where $R=R_{\calT,\s}$ denotes the
ring generated over~$\Z$ by the numbers $(1-z_j)\i$. We write the monomial $T$ as
$c(T) h^n \bx^\bK/\bK!$ with
\be
\label{nAndK}
n = \sum_{j,r,k}\l_j(r,k)\left(r+\frac k2-1\right),\qquad K_j=\sum_{r,k}\l_j(r,k)k
\ee
and where for notational convenience have written $\bx$ and $\bK$ for the
$N$-tuples $(x_1,\dots,x_N)$ and~${(K_1,\dots,K_N)}$ (thus deviating from the
convention in the rest of the paper where boldface denotes matrices) and
$\bx^\bK/\bK!$ for the divided power \smash{$\prod x_j^{K_j}/K_j! $}. To prove the theorem,
we have to bound both the denominators of the numerical coefficient~$c(T)$ and the
further denominators coming from the Gaussian averaging
$ \bx^\bK/\bK! \mapsto\big\langle\bx^\bK/\bK! \big\rangle $.

We begin with the latter question. For this, we recall first that the Gaussian
average $\langle f\rangle_Q$ is given by \smash{${\rm e}^{\D_Q}(f)\bigl|_{x=0}$} for any power
series~$f$, where $\D_Q$ is the Laplacian associated to~$Q$, and hence is equal to
$\D_Q^\ell(f)/\ell!$ if $f$ is a homogeneous polynomial of degree $2\ell$. (For
polynomials of odd degree it of course vanishes trivially.) The Laplacian $\D_Q$
is a quadratic polynomial in the derivatives $\pt_i=\pt/\pt x_i$, and we can enlarge
the ring~$R$ by adjoining to it the coefficients of this polynomial, so since the
image of any divided factorial under any product $\pt_1^{\ell_1}\cdots\pt_N^{\ell_N}$
is an integer, we then certainly have that the Gaussian integral
$\big\langle\bx^\bK/\bK! \big\rangle$ is $1/\ell!$ times an element of~$R$. Unfortunately,
it turns out that this estimate is not good enough for our purposes, and we have
to work a little harder.

Writing $\D_Q$ as an $R$-linear combination of binomials $\pt_i\pt_j$ with
$1\le i\le j\le N$, and applying the multinomial theorem, we see that
$\D_Q^\ell/\ell!$ is an $R$-linear combination of terms
\smash{$\prod_{i\le j}(\pt_i\pt_j)^{m_{ij}}/m_{ij}!$} with $m_{ij}\in\Z_{\ge0}$. Define an even
symmetric $N\times N$ matrix $M$ by setting $M_{ij}=M_{ji}=m_{ij}$ for $i<j$ and
$M_{ii}=2 m_{ii}$. Then \smash{$\prod_{i\le j}(\pt_i\pt_j)^{m_{ij}}=\prod_j\pt_j^{K_j}$}
with $\bK=M\mb1$, where $\mb1$ is the vector consisting of~$N$~1's, and this
sends $\bx^\bK/\bK!$ to~1 and all other monomials to~0.
It follows that a universal denominator of $\big\langle\bx^\bK/\bK! \big\rangle$ is
the number
\begin{align*}
\D(\bK) := \text{l.c.m.}\Biggl\{ \prod_{1\le i\le N}(M_{ii}/2)!\prod_{1\le i<j\le N}M_{ij}!
 \,\Big|\, \text{$M=M^t\in M_{N,N}(\Z_{\ge0})$ even, $M\mb1=\bK$}\Biggr\} .
\end{align*}
Notice that this does divide~$\ell!$, because $\ell$ is the sum of the diagonal
$M_{ii}/2$ and of the $M_{ij}$ with~${i<j}$, so that this statement refines the bound
given above, but $\D(\bK)$ is in general much smaller than~$\ell!$ and this
improvement will be needed for the proof. A usually sharper multiplicative upper
bound for $\D(\bK)$ is the largest integer~$S$ whose square divides the product of
the~$K_j!$ (the proof of this also uses only the integrality of multinomial
coefficients), and then of course~$\D(\bK)$ also divides the g.c.d.\ of $\ell!$ and
$S$, which in general is smaller than either one. (For instance, for $\bK=(6,9,9,10)$
we have $\ell!=355687428096000$, $S=3135283200$, and~${(\ell!,S)=S/3}$.)
Either of these two latter upper bounds would be sufficient for our proof, but in
fact there is an easy upper bound that is stronger than either one of them and is
extremely sharp (in particular, it
is equal to $\D(\bK)$ for all $\bK$ with $N\le4$ and $\max(K_j)\le 30$), namely
\[
\D^*(\bK) := \prod_{\text{$p$ prime}} p^{\delta_p(\bK)} \quad \text{with} \quad
\delta_p(K_1,\dots,K_N) := \sum_{s\ge1}
\left[\frac12\sum_{j=1}^N\left[\frac{K_j}{p^s}\right]\right] .
\]
To show that $\D(\bK)$ divides~$\D^*(\bK)$, we need
$\sum_iV_p(M_{ii}/2)+\sum_{i<j}V_p(M_{ij})\le\delta_p(\bK)$ for every prime~$p$ and
every even symmetric matrix~$M$ with non-negative entries and row sums~$\bK$, where~${V_p(m):=v_p(m!)}$ for any~$m\ge0$ denotes the largest power of~$p$ dividing~$m!$.
In view of the standard formula $V_p(m)=\sum_{s\ge1}[m/p^s]$, it suffices for this
to show that
\[\sum_i[M_{ii}/2q]+\sum_{i<j}[M_{ij}/q]\le\frac12\sum_j[K_j/q]\]
for every prime power~$q$, and this follows immediately from the obvious facts
$[x/2q]=[[x/2]/q]$ and $[x]+[y]+\cdots \le[x+y+\cdots]$ valid for arbitrary real
numbers $x, y,\dots$.

Now going back to our main problem, we now see that it suffices to show that the
product~$\D^*(\bK) c(T)$ has denominator dividing~$D_n$ for all terms~$T$ as above,
with $\bK=(K_1,\dots,K_N)$ and~$n$ defined by~\eqref{nAndK}. We will prove this
one prime at a time (ignoring the prime~2), which is natural in view of the fact
that the upper bound~$\D^*(\bK)$ is defined by its prime power decomposition. To do
this, we will split both the term~$T$ and the corresponding numerical
coefficient~$c(T)$, and also each of the $N$ factors $T_j$ and $c(T_j)$ of which
they are comprised, as the product of four factors in a way depending on the
prime~$p$ being studied, labelled ``$s$'' (terms with $r=0$ and~$k$ smaller than~$p$),
``$b$'' (terms with $r=0$ and $k$ bigger than or equal to~$p$),
``1''~(terms with~${r=1}$), and ``$\ge2$'' (terms with $r\ge2$), with a similar splitting
of the individual weights~$K_j$ into the sum of four pieces
\smash{$K_j^{(s)}=\sum_{3\le k<p}\l_j(0,k) k$}, \smash{$K_j^{(b)}=\sum_{k\ge p}\l_j(0,k) k$}, \smash{$
K_j^{(1)}=\sum_{k\ge1}\l_j(1,k) k$}, and~\smash{$K_j^{(\ge2)}=\sum_{r\ge2, k\ge0}\l_j(r,k) k$}.
We also decompose the number~$n$ (the exponent of~$h$ in~$T$) in~\eqref{nAndK} as
$\ell+n'-t$ with
\[
\ell := \frac12 \sum_{j=1}^NK_j , \qquad
n' :=\sum_{\substack{1\le j\le N \\ r\ge2, k\ge0\substack}}\l_j(r,k) (r-1) , \qquad
t :=\sum_{\substack{1\le j\le N \\k\ge3\substack}}\l_j(0,k)
\]
and also split $t$ as \smash{$t^{(s)}+t^{(b)}$} according as $3\le k\le p-1$ or $k\ge p$ in
the summation, with each of \smash{$t^{(s)}$} and \smash{$t^{(b)}$} splitting into the sum over
$1\le j\le N$ of pieces \smash{$t_j^{(s)}$} and \smash{$t_j^{(b)}$} in the obvious way.

The numerical coefficient $c(T)$ can be decomposed as
\be
\label{cTdecomp}
c(T) = \prod_{1\le j\le N}\frac{K_j!}{P_j(T)} \cdot
\prod_{\substack{1\le j\le N \\ r\ge2, k\ge0\substack}}
\frac1{\l_j(r,k)!} \biggl(\frac{B_r}{r!}\biggr)^{\l_j(r,k)}
\ee
with
\[
P_j(T) = \prod_{0\le r\le1, k\ge1}\l_j(r,k)! k!^{\l_j(r,k)} \cdot
\prod_{r\ge2, k\ge0}k!^{\l_j(r,k)} ,
\]
which we can split up further as \smash{$P_j^{(s)}(T)P_j^{(b)}(T)P_j^{(1)}(T)P_j^{(\ge2)}(T)$}.
The reason that we have included the factor $\l_j(r,k)$ into the definition of
$P_j(T)$ for $0\le r\le1$ but not for~$r\ge2$ is that~$\l!k!^\l$ divides $(k\l)!$
for all $\l\ge0$ if $k$ is strictly positive but not if $k=0$, and the terms with~${r\ge2}$ can have~$k=0$. Then the product \smash{$P_j^{(b)}(T)P_j^{(1)}(T)P_j^{(\ge2)}(T)$}
divides \smash{$\big(K_j-K_j^{(s)}\big)!$}, while the first factor \smash{$P_j^{(s)}(T)$} divides \smash{$t_j^{(s)}!$}
up to a $p$-adic unit because here $k$ is always less than~$p$ and therefore~$k!$
is not divisible by~$p$. (Here we have made repeated use of the integrality of
multinomial coefficients.) On the other hand, by Lemma~\ref{LemmaB} below
and the submultiplicativity of~$D_n$, the second factor in~\eqref{cTdecomp} has
denominator dividing~$D_{n'}$. Putting this all together, we deduce that~$c(T)$
is $G(T)/D_{n'}$ times a $p$-adic integer for every~$p$ (always different from~2
and not dividing the denominators of the elements of~$R$), where
\smash{$G(T)=\prod_{j=1}^N\bigl(K_j!/t_j^{(s)}! \big(K_j-K_j^{(s)}\big)!\bigr)$}. Using the
submultiplicativity of $D_n$ again, this reduces the problem to showing that
${\delta_p(\bK)\le v_p(D_{\ell-t} G(T))}$ for each~$p$, and in view of
the definitions of $\delta_p(\bK)$ and $D_n$ and of the above-mentioned formula~${V_p(m)=\sum_{s\ge1}[m/p^s]}$ for the $p$-adic valuation of factorials, this in turn
will follow if we can show that
\be
\label{Final}
\Biggl[\frac12\sum_{j=1}^N\biggl[\frac{K_j}q\biggr]\Biggr] \le
\biggl[\frac{\ell-t}{q^*}\biggr]
+ \sum_{j=1}^N \Bigg(\biggl[\frac{K_j}q\biggr] \m \Bigg[\frac{t_j^{(s)}}q\Bigg]
\m \Bigg[\frac{K_j-K_j^{(s)}}q\Bigg] \Bigg)
\ee
for each prime power $q=p^s$ with $s\ge1$, where $q^*:=p^{s-1}(p-2)$. % =(1-2/p) q$.

For this final step, we first note that
\[
\frac{\ell-t}{q^*} = \sum_{j=1}^N
\frac{K_j^{(s)}+K_j^{(b)}+K_j^{(1)}+K_j^{(\ge2)}-2t_j^{(s)}-2t_j^{(b)}}{2q^*}
 \ge \sum_{j=1}^N\frac{K_j-2t_j^{(s)}}{2q}
\]
since
\[
\frac{\big(K_j^{(b)}-2t_j^{(b)}\big)}{q^*}\ge \left(1-\frac{2}{p}\right) \frac{K_j^{(b)}}{q^*}=\frac{K_j^{(b)}}{q}
\]
 \big(because $k\ge p$ in the terms defining \smash{$K_j^{(b)}$} and \smash{$t_j^{(b)}$}\big) and~$q^*<q$.
Using that $\big[\frac x{2q}\big]=\big[\frac12\big[\frac xq\big]\big]$, we deduce that
\[
\biggl[\frac{\ell-t}{q^*}\biggr] \ge
\Biggl[\frac12\sum_{j=1}^N\Bigg[\frac{K_j-2t_j^{(s)}}{q}\Bigg]\Biggr]
\]
and hence (since $x\le y$ certainly implies $[x/2]\le[y/2]$) the
inequality~\eqref{Final} will follow if we have the inequality
\[
\biggl[\frac{K_j}q\biggr] \le \Bigg[\frac{K_j-2t_j^{(s)}}{q}\Bigg]
+ 2 \Bigg(\biggl[\frac{K_j}q\biggr] \m \Bigg[\frac{t_j^{(s)}}q\Bigg]
\m \Bigg[\frac{K_j-K_j^{(s)}}q\Bigg] \Bigg)
\]
for every $1\le j\le N$. But this inequality is trivial, since
\smash{$K_j-K_j^{(s)}\le K_j-3t_j^{(s)}\le K_j-2t_j^{(s)}$} (because every $k$ in the
definition of \smash{$K_j^{(s)}$} is $\ge3$) and
\[
\biggl[\frac{K_j}q\biggr]\ge\Bigg[\frac{K_j-2t_j^{(s)}}q\Bigg]
+2\Bigg[\frac{t_j^{(s)}}q\Bigg].
\]
 This completes the proof of Theorem~\ref{thm.denom} modulo that of the following lemma.
\end{proof}

\begin{Lemma}\label{LemmaB}
For any integers $r\ge2$ and $\l\ge0$, we have
\be
\label{BernEstimate}
\frac1{\l!} \left(\frac{B_r}{r!}\right)^\l \in \frac1{D_{\l(r-1)}} \Z .
\ee
\end{Lemma}

\begin{proof}
We prove this one prime at time. By well-known results of
von Staudt and Clausen, the Bernoulli number $B_r$ ($r>0$) has $p$-adic valuation
$-1$ if~${(p-1)|r}$ and $B_r/r$ is $p$-integral if~${p-1}$ does not divide~$r$. From this
we deduce that the $p$-adic valuation of the denominator of $B_r/r!$ is bounded
above by $\bigl[\frac r{p-1}\bigr]$, which is $\le\bigl[\frac {r-1}{p-2}\bigr]$
since $\frac{r-1}{p-2}\ge\frac r{p-1}$ if $r\ge p-1$ and both expressions
vanish otherwise. The $p$-adic valuation of the denominator of
\smash{$\frac1{\l!}\bigl(\frac{B_r}{r!}\bigr)^\l$} is therefore bounded above by
\smash{$V_p(\l)+\l\bigl[\frac{r-1}{p-2}\bigr]$}. On the other hand, from the definition
of~$D_n$ we have $v_p(D_n)=m+V_p(m)=V_p(pm)$, where $m=\bigl[\frac n{p-2}\bigr] $.
We must therefore show that
\[
V_p(\l)+\l \biggl[\frac{r-1}{p-2}\biggr] \le
V_p\biggl(p \biggr[\frac{\l(r-1)}{p-2}\biggr]\biggl)
\]
for all $\l\ge0$ and $r\ge2$. For this, we make a case distinction: if $ 2\le r<p-1$, then
\[
\text{l.h.s.} = V_p(\l) = V_p\biggl(p \biggl[\frac\l p\biggr]\biggr) \le
V_p\biggl(p \biggl[\frac\l{p-2}\biggr]\biggr) \le \text{r.h.s.} ,
\]
while if $r\ge p-1$, then we set $h=\bigl[\frac{r-1}{p-2}\bigr]\ge1$ and have
instead
\[
\text{l.h.s.} = V_p(\l)+\l h \le \l h+V_p(\l h) = V_p(p\l h) \le \text{r.h.s.}
\]
because $[\l x]\ge \l [x]$ for any positive real number~$x$.
\end{proof}
\begin{thebibliography}{99}
\bibitem[10]{CGZ}
Calegari F., Garoufalidis S., Zagier D., Bloch groups, algebraic {$K$}-theory,
 units, and {N}ahm's conjecture, \href{https://doi.org/10.24033/asens.2537}{\textit{Ann. Sci. \'Ec. Norm. Sup\'er.~(4)}}
 \textbf{56} (2023), 383--426, \href{https://arxiv.org/abs/1712.04887}{arXiv:1712.04887}.

\bibitem[14]{DG}
Dimofte T., Garoufalidis S., The quantum content of the gluing equations,
  \href{https://doi.org/10.2140/gt.2013.17.1253}{\textit{Geom. Topol.}} \textbf{17} (2013), 1253--1315, \href{https://arxiv.org/abs/1202.6268}{arXiv:1202.6268}.

\bibitem[15]{DG2}
Dimofte T., Garoufalidis S., Quantum modularity and complex {C}hern--{S}imons
 theory, \href{https://doi.org/10.4310/CNTP.2018.v12.n1.a1}{\textit{Commun. Number Theory Phys.}} \textbf{12} (2018), 1--52,
 \href{https://arxiv.org/abs/1511.05628}{arXiv:1511.05628}.

\bibitem[70]{OEIS}
Sloane N., Online encyclopedia of integer sequences, available at
 \url{https://oeis.org}.

\end{thebibliography}
\end{document}
